\section{Introduction}
\label{sec:intro}

Medical image segmentation models are trained on one acquisition setting and
routinely deployed on another: a different scanner, sequence, or imaging
modality than the one seen during training. Because MRI contrast is a
protocol choice rather than a physical constant, and because CT, ultrasound,
and MRI encode the same anatomy in incompatible intensity distributions, a
segmentation network that has only ever seen T1-weighted brain images can
fail outright on T2-weighted or FLAIR images of the same anatomy, let alone a
different modality. Collecting and annotating training data for every
contrast and modality a model might encounter in deployment is not
practical: manual segmentation is expensive, and clinical sites rarely share
raw data across institutions. The single-source domain generalization (DG)
setting --- train on one labelled dataset in one contrast, generalize to
unseen contrasts and modalities without any target-domain data --- is
therefore the setting that matters for real deployment, and it is the
setting we address.

The dominant single-source strategy is label-generative domain randomization,
popularized by SynthSeg~\cite{billot2023synthseg}: given a dense anatomical
label map, intensities are resampled from a per-label Gaussian mixture,
producing images of arbitrary, unrealistic contrast that force the network to
ignore contrast and rely on shape and layout. This is effective, but it comes
with two structural costs. First, it is \emph{label-driven}: it requires a
dense anatomical segmentation of every training subject to resample from, so
generating synthetic contrasts costs exactly the expensive annotation the
method is trying to reduce reliance on. Second, and less discussed, it is
\emph{texture-destroying}: replacing every labelled region with piecewise
Gaussian noise erases the real gradients, partial-volume boundaries, and
unlabelled fine structure that the source image contains. Layout is
preserved; texture is not.

We ask whether texture-destroying synthesis is necessary for contrast
robustness, or merely the path of least resistance once a label map is
assumed. We introduce \textbf{PALETTE} (\emph{Partition-based Affine
Local-intEnsity Transform for Texture-prEserving augmentation}), a
single-source augmentation that requires no label map at all. PALETTE
partitions the foreground of a scan into intensity classes with unsupervised
1-D $k$-means, further subdivides each class into spatial sub-regions with a
Voronoi tessellation, and applies a signed random affine remap $y = \mu +
\alpha(x-\bar{x})$, $\alpha \in \pm U(0.5, 2.0)$, to the \emph{real}
intensities within each sub-region. Because the transform remaps existing
intensities rather than replacing them with sampled noise, edges, gradients,
and partial-volume detail are preserved within each sub-region; only the
region-to-region contrast relationship is randomized, including sign
inversions that no physical acquisition protocol produces. The result is a
label-free augmentation that manufactures large, systematic contrast
diversity while leaving the texture cues a segmentation network can exploit
intact.

This distinction --- image-driven, texture-preserving remapping versus
label-driven, texture-destroying resampling --- is not a matter of degree:
the two families make opposite bets about what a network should be invariant
to. We test this bet directly. Training a single nnU-Net~\cite{isensee2021nnunet}
configuration with PALETTE, tuned once and reused unchanged across every
dataset, we evaluate on eight cross-contrast and cross-modality settings
spanning abdominal CT/MRI (CHAOS $\to$ AMOS/SLIVER07), brain MRI
(OpenHarmony, OpenMS), and brain tumor MRI (BraTS), against a
no-augmentation baseline and both variants of SynthSeg (the original,
label-generative noEM variant, and a label-free EM variant for which no
public implementation exists and which we re-implement as a fair baseline in
its own right). PALETTE outperforms the no-augmentation baseline and
SynthSeg-noEM significantly in all eight settings, and beats the
stronger SynthSeg-EM baseline significantly in seven of eight. The margin on
aggregate whole-organ Dice is often thin --- the setting is genuinely hard,
and the strongest competing configuration is our own ablation without the
PALETTE transform --- but the margin is large precisely where the texture
argument predicts it should be: on fine, texture-defined structures and on
the most contrast-shifted test modalities, where destroying real texture
costs the most.

A thin average is not evidence for a mechanism; a mechanism claim needs a
controlled test. We isolate texture preservation as the causal factor,
rather than an incidental correlate of PALETTE's partitioning scheme, with a
causal ablation: we hold the $k$-means/Voronoi partition and the per-region
$(\mu,\alpha)$ statistics fixed, and replace only the real-intensity affine
remap with label-style noise fill, matching SynthSeg's resampling within
PALETTE's own geometry. Structures that are texture-defined collapse under
this ablation while boundary-defined structures do not, isolating texture
preservation --- not the partitioning scheme it happens to share with
SynthSeg --- as the causal source of PALETTE's advantage.

In summary, we make the following contributions:
\begin{itemize}
    \item \textbf{PALETTE}, a label-free, single-source contrast augmentation
    that partitions a scan by intensity and space and applies a signed affine
    remap of the real intensities per sub-region, preserving texture and
    gradients that label-generative synthesis destroys.
    \item A systematic evaluation across eight cross-contrast and
    cross-modality segmentation settings (abdominal CT/MRI, brain MRI, brain
    tumor MRI) with a single, untuned-per-dataset configuration, against a
    no-augmentation baseline and both variants of SynthSeg, including a
    re-implementation of the label-free SynthSeg-EM variant for which no
    public code exists.
    \item A causal ablation isolating texture preservation, rather than the
    shared intensity/spatial partitioning scheme, as the source of PALETTE's
    advantage on texture-defined structures.
\end{itemize}
